\documentclass[11pt]{amsart}
\usepackage{amsmath,amssymb,amsthm,booktabs}
\usepackage[margin=1.15in]{geometry}
\usepackage[hidelinks]{hyperref}
\newtheorem{theorem}{Theorem}[section]
\newtheorem{lemma}[theorem]{Lemma}
\theoremstyle{remark}
\newtheorem{remark}[theorem]{Remark}
\newcommand{\R}{\mathbb{R}}
\newcommand{\Z}{\mathbb{Z}}
\newcommand{\PP}{\mathbb{P}}
\newcommand{\E}{\mathbb{E}}
\newcommand{\Var}{\operatorname{Var}}
\newcommand{\dimH}{\dim_{\mathrm H}}
\newcommand{\dimB}{\overline{\dim}_{\mathrm B}}
\newcommand{\doi}[1]{\href{https://doi.org/#1}{doi:\nolinkurl{#1}}}
\setlength{\emergencystretch}{2.5em}

\title[Double times on slices of the Brownian sheet]{Double times on slices of the Brownian sheet in dimensions four and five}
\author{Yitzchak Shmalo}
\address{Einstein Institute of Mathematics, The Hebrew University of Jerusalem, Jerusalem, Israel}
\email{yspennstate@gmail.com}
\date{November 2026}
\subjclass[2020]{60G60, 60G15, 60J65, 28A80}
\keywords{Brownian sheet, slices, double points, Hausdorff dimension, box dimension, exceptional times}

\raggedbottom
\begin{document}
\begin{abstract}
Let $B$ be a $d$-dimensional Brownian sheet. For $s>0$ the slice $t\mapsto B(s,t)$ is a multiple of $d$-dimensional
Brownian motion, so for $d\ge4$ a fixed slice has no double points, while for $d\in\{4,5\}$ there are random exceptional
slices that have them. Khoshnevisan asked whether, outside a single null set, the set of double times of every slice has
Hausdorff dimension zero when $d\in\{4,5\}$. We prove that it does: almost surely, for every $s>0$ the set of pairs
$u\ne v$ with $B(s,u)=B(s,v)$ has Hausdorff dimension zero, and its part in a product of two disjoint compact intervals
has upper box dimension zero. For $d=5$ both statements also follow from the finiteness of these sets stated by Chen;
the new case is $d=4$. The proof bounds high moments of the number of near-coincidences on a grid, using a
conditional variance estimate obtained from Brownian bridges, and then takes a union over a grid of slices whose cost is
absorbed by the choice of the moment.
\end{abstract}
\maketitle

\section{Introduction}\label{sec:intro}

Let $B=(B_1,\dots,B_d)$ be a continuous $\R^d$-valued Brownian sheet on $[0,\infty)^2$: the coordinates $B_i$ are
independent centred Gaussian processes with
\[
\E[B_i(s,t)B_i(s',t')]=\min(s,s')\min(t,t').
\]
For fixed $s>0$ the process $t\mapsto s^{-1/2}B(s,t)$ is a standard Brownian motion in $\R^d$, which has no double
points when $d\ge4$. Some slices do have double points. Let
\[
D_d=\{s\ge0:\ B(s,u)=B(s,v)\text{ for some }0<u<v\}.
\]
Lyons~\cite{Lyo86} proved that $D_d\ne\{0\}$ with positive probability if and only if $d\le5$, and Mountford~\cite{Mou90}
proved that $\dimH D_d=\min(1,(3-d/2)_+)$ almost surely; see~\cite[(1.10) and (1.12)]{Kho05}. For $d\le3$ every slice
has double points~\cite{Pen89,Pen90}. So $d\in\{4,5\}$ is the range in which the exceptional slices are rare (in
$d=4$ they form a set of dimension one and Lebesgue measure zero) but present.

For $s\ge0$ Khoshnevisan~\cite[(8.2)]{Kho05} considers
\[
Y_d(s)=\{(t_1,t_2)\in[0,\infty)^2:\ B(s,t_1)=B(s,t_2)\},
\]
and asks, as Open Problem~2 of~\cite[Section~8.2]{Kho05}, whether for $d\in\{4,5\}$ one has, outside a single null set,
$\dimH Y_d(s)=0$ for all $s\ge0$. He conjectures that the answer is yes. As written, $Y_d(s)$ contains the diagonal
$\{(t,t):t\ge0\}$, which has dimension one, and $Y_d(0)=[0,\infty)^2$ because $B(0,\cdot)=0$; the question concerns the
double times off the diagonal on the slices $s>0$. For $s>0$ let
\begin{equation}\label{eq:Ystar}
Y^*_d(s)=\{(u,v)\in[0,\infty)^2:\ u\ne v,\ B(s,u)=B(s,v)\}.
\end{equation}
We write $\dimB$ for the upper box dimension, and we use the conventions $\dimH\emptyset=\dimB\emptyset=0$.

\begin{theorem}\label{thm:main}
Let $d\in\{4,5\}$. Almost surely the following holds simultaneously for all $s>0$.
\begin{enumerate}
\item[(a)] $\dimH Y^*_d(s)=0$.
\item[(b)] For all compact intervals $I,J\subset[0,\infty)$ with $I\cap J=\emptyset$,
$\dimB\bigl(Y^*_d(s)\cap(I\times J)\bigr)=0$.
\end{enumerate}
\end{theorem}

Consequently, outside a single null set, $\dimH Y_d(s)=1$ for every $s>0$, the dimension coming from the diagonal alone.

\begin{remark}\label{rem:scope}
Theorem~\ref{thm:main} bounds the size of the double-time sets, not their number of points. The bound is compatible with
nonempty and with infinite sets $Y^*_4(s)$, and with $D_4$ having dimension one. Khoshnevisan's Open Problem~3 asks
whether, for $d\in\{4,5\}$, outside a single null set $Y_d(s)$ is finite for all $s>0$~\cite[Section~8.2]{Kho05}.
Chen~\cite[Theorems~1.3.3 and~1.3.4]{Che26} states that the off-diagonal double-time sets are finite for every $s>0$
when $d=5$, and that in $d=4$ every interval of slice parameters contains slices with infinitely many double times. A
finite set has Hausdorff and upper box dimension zero, so for $d=5$ both parts of Theorem~\ref{thm:main} follow from
the first statement; the case $d=4$ is new, and the projection argument of Section~\ref{sec:proof} gives $d=5$
without~\cite{Che26}. At~\cite[(4.42)]{Che26} the field $F(x,y)=V(y)-U(x)$, built from two independent increments
$U$ and $V$ of a slice, is treated as a Brownian sheet. Its covariance is $s(\min(x,x')+\min(y,y'))$, a sum rather
than a product, and its rectangular increments vanish, so that step of the printed proof needs another argument. The
proof below does not use these results.
\end{remark}

\begin{remark}\label{rem:low}
For $d\le3$ the answer is no. At a fixed $s>0$ the slice is a multiple of Brownian motion in $\R^d$, whose set of double
times $\{(u,v):u<v,\ W(u)=W(v)\}$ has Hausdorff dimension $2-d/2>0$ almost surely~\cite{Ros84}; see
also~\cite[(4.2)]{KWX06}. This is the case that Khoshnevisan connects with intersection local
times~\cite[Section~8.2]{Kho05}.
\end{remark}

\textbf{Why counting dimensions is not enough.} Away from the diagonal and from $s=0$, the field
$X(s,u,v)=B(s,v)-B(s,u)$ has three parameters, each with H\"older exponent just below $1/2$, and $d$ independent
coordinates. The zero set of $X$ has dimension $3-d/2$ for $d<6$~\cite[Theorem~5.1 and (8.1)]{Kho05}, and its
projection on the $s$-axis, the set $D_d$, has dimension $\min(1,3-d/2)$ by~\cite[(1.12)]{Kho05}. Subtracting
suggests that the fibres have dimension $\max(0,2-d/2)$, which is $0$ for $d\in\{4,5\}$. Dimension does not subtract
fibre by fibre: the union of a horizontal and a vertical segment in the plane has dimension one, its projection on the
first axis has dimension one, and one of its fibres is a segment. Theorem~\ref{thm:main} is a statement about every
fibre, and its proof has to control all slices at once.

\textbf{The idea of the proof.} Fix a compact interval $K=[a,b]$ of slice parameters with $a>0$ and two disjoint
compact time intervals $I$ and $J$. For a grid of mesh $\delta=2^{-n}$ and a radius $r$ with $\delta\le r^2\le1$, let
$N_n(s,r)$ be the number of grid pairs $(u,v)\in I\times J$ with $|B(s,v)-B(s,u)|\le r$. For fixed $s$ we show
\[
\E\bigl[N_n(s,r)^p\bigr]\le C_p\bigl(r^4\delta^{-2}(1+\log(1/\delta))\bigr)^p
\qquad\text{for every }p\ge1
\]
in $d=4$ (Lemma~\ref{lem:moments}). The point is that the exponent of $r$ matches the exponent of $\delta$: the two time
coordinates contribute $\delta^{-2}$, the four spatial coordinates contribute $r^4$, and their critical balance leaves
only a logarithm. The moment bound comes from a lower bound on the variance of $B(s,v)-B(s,u)$ conditionally on earlier
pair differences, which we obtain by conditioning further on the Brownian motion $t\mapsto B(s,t)$ at the endpoints of
the earlier pairs and using that the remaining randomness consists of independent Brownian bridges
(Lemma~\ref{lem:var}). With $r$ of order $\delta^{\alpha}$ for $\alpha$ close to $1/2$, Markov's inequality makes
$N_n(s,r)>\delta^{-\eta}$ very unlikely, so unlikely that a union over the slices of the grid, at most a constant multiple of $\delta^{-1}$ of them, is
affordable once $p$ is large. H\"older continuity of $B$ then transfers the bound to every real $s\in K$ and every exact
double time.

The proof uses only the Gaussian structure of a single slice and the H\"older continuity of the sheet. The sectorial
local nondeterminism of Khoshnevisan, Wu and Xiao~\cite[Lemma~1.1]{KWX06} and their moment estimates for clusters of
nearby values~\cite[Lemma~3.1]{KWX06} play a similar role for the full sheet; for a slice and two separated time
intervals the Brownian bridge bound below is elementary. The double points of the full sheet, a different question, are
treated in~\cite{DKNWX12,DM15}. Chen and Woessner~\cite{CW25} prove a statement that holds simultaneously over a
continuum of parameters for the zero set of the sheet: almost surely no projection of large enough rank changes its
Hausdorff dimension, which answers another question of Khoshnevisan.

Section~\ref{sec:var} proves the conditional variance and small-ball bounds, Section~\ref{sec:moments} the moment bound,
and Section~\ref{sec:proof} the theorem. Throughout, $C$ denotes a finite constant whose dependence is indicated, and
whose value may change from line to line.

\section{Conditional variance and small balls}\label{sec:var}

Fix $0<a<b$, $K=[a,b]$, and compact intervals $I=[A_1,A_2]$ and $J=[A_3,A_4]$ with $0\le A_1<A_2<A_3<A_4$. Fix $s\in K$.
For pairs $z_j=(u_j,v_j)\in I\times J$, $j=1,\dots,k$, put
\[
X_j=B(s,v_j)-B(s,u_j)\in\R^d.
\]
The pairs need not be distinct, and their coordinates may coincide with each other or with the endpoints of $I$ and
$J$. For a further pair $(u,v)\in I\times J$ put $X=B(s,v)-B(s,u)$, and let
\[
U=\{A_1,A_2,u_1,\dots,u_k\},\qquad V=\{A_3,A_4,v_1,\dots,v_k\},
\]
\[
\rho_U(u)=\min_{x\in U}|u-x|,\qquad \rho_V(v)=\min_{y\in V}|v-y|,\qquad \rho=\rho_U(u)+\rho_V(v).
\]
We call the points of $U\cup V$ anchors.

\begin{lemma}\label{lem:var}
For each coordinate $i$,
\begin{equation}\label{eq:var}
\Var\bigl(X^{(i)}\bigm| X^{(i)}_1,\dots,X^{(i)}_k\bigr)\ge\frac a2\,\rho,
\end{equation}
where $X^{(i)}$ and $X^{(i)}_j$ are the $i$-th coordinates of $X$ and $X_j$.
\end{lemma}

\begin{proof}
The process $W(t)=B_i(s,t)$ is a Brownian motion with variance rate $s$. For a centred Gaussian family the conditional
variance of one member $Z$ given others is the squared $L^2$ distance from $Z$ to their linear span; it is
deterministic and does not increase when the family is enlarged. Let $\mathcal F$ be generated by $X^{(i)}_1,\dots,X^{(i)}_k$ and $\mathcal G$
by the values of $W$ at all anchors. Each $X^{(i)}_j=W(v_j)-W(u_j)$ is $\mathcal G$-measurable, so $\mathcal F\subset\mathcal
G$, and $\Var(X^{(i)}\mid\mathcal F)\ge\Var(X^{(i)}\mid\mathcal G)$.

On an interval $[\ell,m]$ between consecutive anchors let $\beta(t)=W(t)-\bigl((m-t)W(\ell)+(t-\ell)W(m)\bigr)/(m-\ell)$.
These residuals are centred Brownian bridges, independent of $\mathcal G$ and of the residuals on the other intervals,
and $W(t)-\beta(t)$ is $\mathcal G$-measurable, so the conditional variance of $W(t)$ given $\mathcal G$ is that of
$\beta(t)$. Since $A_1,A_2\in U$ and every point of $V$ exceeds $A_2$, the
anchors adjacent to $u$ lie in $U$. If $u\in U$ its conditional variance is $0=\rho_U(u)$. Otherwise let $\ell<u<m$ be
the adjacent anchors; the variance of the bridge at $u$ is
\[
s\,\frac{(u-\ell)(m-u)}{m-\ell}\ge\frac s2\min(u-\ell,\,m-u)=\frac s2\,\rho_U(u).
\]
The same holds for $v$ with $V$. The points $u$ and $v$ lie in different intervals between anchors, because $A_2$ and
$A_3$ are anchors, so their bridges are independent and
$\Var(X^{(i)}\mid\mathcal G)\ge\frac s2(\rho_U(u)+\rho_V(v))\ge\frac a2\rho$. Combined with the first paragraph this
proves~\eqref{eq:var}. No invertibility of the covariance matrix of $X^{(i)}_1,\dots,X^{(i)}_k$ is needed.
\end{proof}

\begin{lemma}\label{lem:smallball}
Let $d=4$. There is a constant $C$ depending only on $a$ such that for every $r>0$,
\begin{equation}\label{eq:smallball}
\PP\bigl(|X|\le r\bigm| X_1,\dots,X_k\bigr)\le\frac{C\,r^4}{(r^2+\rho)^2}
\end{equation}
almost surely, whatever the values of $X_1,\dots,X_k$.
\end{lemma}

\begin{proof}
The coordinates of $B$ are independent, so given $X_1,\dots,X_k$ the coordinates of $X$ are independent Gaussian
variables; the $i$-th has the conditional law of $X^{(i)}$ given $X^{(i)}_1,\dots,X^{(i)}_k$, with some mean $\mu_i$
and the variance $\sigma^2$ of Lemma~\ref{lem:var}, which is the same for all $i$. If $\rho>0$ then $\sigma^2\ge
a\rho/2>0$ and the conditional density of $X$ is at most $(2\pi\sigma^2)^{-2}$, so
\[
\PP\bigl(|X|\le r\bigm| X_1,\dots,X_k\bigr)\le\frac{\pi^2r^4}{2}\,(2\pi\sigma^2)^{-2}=\frac{r^4}{8\sigma^4}\le
\frac{r^4}{2a^2\rho^2}.
\]
If $\rho\le r^2$ the right side of~\eqref{eq:smallball} is at least $C/4$, which is at least $1$ for $C\ge4$. If
$\rho>r^2$ then $(r^2+\rho)^2\le4\rho^2$ and the display gives the bound with $C=2/a^2$. Hence~\eqref{eq:smallball}
holds with $C=\max(4,2/a^2)$, including the case $\rho=0$, in which the conditional law may be degenerate.
\end{proof}

\section{Moments of the number of near-coincidences}\label{sec:moments}

Let $\delta=2^{-n}$ and assume that $A_1,\dots,A_4$ lie in $\delta\Z$. Write $I_\delta=I\cap\delta\Z$ and
$J_\delta=J\cap\delta\Z$.

\begin{lemma}\label{lem:lattice}
Let $0<\delta\le r^2\le1$, and let $U\subset I$ and $V\subset J$ be nonempty finite sets with at most $k+2$ elements each. Then
\begin{equation}\label{eq:lattice}
\sum_{u\in I_\delta}\sum_{v\in J_\delta}\frac{r^4}{(r^2+\rho_U(u)+\rho_V(v))^2}\le C\,(k+2)^2\,r^4\delta^{-2}
\bigl(1+\log(1/\delta)\bigr),
\end{equation}
where $C$ depends only on the lengths of $I$ and $J$.
\end{lemma}

\begin{proof}
Assign each $u\in I_\delta$ to a nearest point $x\in U$, so that $\rho_U(u)=|u-x|$, and each $v\in J_\delta$ to a
nearest point $y\in V$. For fixed $x$ and an integer $j\ge0$, each of the half-open intervals $[x+j\delta,x+(j+1)\delta)$
and $(x-(j+1)\delta,x-j\delta]$ contains exactly one point of $\delta\Z$, so at most two points $u$ satisfy
$j\delta\le|u-x|<(j+1)\delta$; the same holds for $y$. Let $L$ be the sum of the lengths of $I$ and $J$. Summing over
the at most $(k+2)^2$ pairs $(x,y)$ and over $j,l\ge0$ with $j+l\le L/\delta$, and putting $q=j+l$, which arises from
$q+1$ pairs $(j,l)$, the left side of~\eqref{eq:lattice} is at most
\[
4(k+2)^2\sum_{q=0}^{\lfloor L/\delta\rfloor}\frac{(q+1)\,r^4}{(r^2+q\delta)^2}
=4(k+2)^2\,r^4\delta^{-2}\sum_{q=0}^{\lfloor L/\delta\rfloor}\frac{q+1}{(r^2/\delta+q)^2}.
\]
Since $r^2/\delta\ge1$, each term of the last sum is at most $1/(q+1)$, and the sum is at most
$1+\log(1+L/\delta)\le C_L(1+\log(1/\delta))$.
\end{proof}

For $s\in K$ and $r>0$ define the number of near-coincidences
\[
N_n(s,r)=\#\bigl\{(u,v)\in I_\delta\times J_\delta:\ |B(s,v)-B(s,u)|\le r\bigr\}.
\]

\begin{lemma}\label{lem:moments}
Let $d=4$. For every integer $p\ge1$ there is a constant $C_p$, depending on $p$, $a$ and the lengths of $I$ and $J$,
such that for all $s\in K$ and all $r$ with $\delta\le r^2\le1$,
\begin{equation}\label{eq:moments}
\E\bigl[N_n(s,r)^p\bigr]\le C_p\Bigl(r^4\delta^{-2}\bigl(1+\log(1/\delta)\bigr)\Bigr)^p.
\end{equation}
\end{lemma}

\begin{proof}
Write $\Lambda=r^4\delta^{-2}(1+\log(1/\delta))$ and $E(z)=\{|B(s,v)-B(s,u)|\le r\}$ for $z=(u,v)$. Then
\[
\E\bigl[N_n(s,r)^p\bigr]=\sum_{z_1,\dots,z_p}\PP\bigl(E(z_1)\cap\dots\cap E(z_p)\bigr),
\]
the sum running over ordered $p$-tuples of pairs in $I_\delta\times J_\delta$, repetitions allowed. Fix
$z_1,\dots,z_{p-1}$, let $U$ and $V$ be the anchor sets built from them as in Section~\ref{sec:var}, and condition on
$X_1,\dots,X_{p-1}$. By Lemma~\ref{lem:smallball},
\[
\PP\bigl(E(z_1)\cap\dots\cap E(z_p)\bigr)\le\PP\bigl(E(z_1)\cap\dots\cap E(z_{p-1})\bigr)\,
\frac{C\,r^4}{(r^2+\rho_U(u_p)+\rho_V(v_p))^2},
\]
because the bound of Lemma~\ref{lem:smallball} does not depend on the values of $X_1,\dots,X_{p-1}$. Summing over $z_p$
with Lemma~\ref{lem:lattice}, where $U$ and $V$ have at most $p+1$ elements, gives the factor $C(p+1)^2\Lambda$, and
leaves the sum over $(p-1)$-tuples. Repeating the step $p$ times gives $\E[N_n(s,r)^p]\le C^p((p+1)!)^2\Lambda^p$. For
$p=1$ the anchor sets are $\{A_1,A_2\}$ and $\{A_3,A_4\}$ and the same argument applies.
\end{proof}

\section{Proof of Theorem~\ref{thm:main}}\label{sec:proof}

\begin{lemma}\label{lem:holder}
Let $R\subset[0,\infty)^2$ be a compact rectangle and $0<\alpha<1/2$. Almost surely there is a finite random $H$ such
that
\begin{equation}\label{eq:holder}
|B(s,t)-B(s',t')|\le H\max(|s-s'|,|t-t'|)^\alpha\qquad\text{for all }(s,t),(s',t')\in R.
\end{equation}
\end{lemma}

\begin{proof}
On $R$, each coordinate satisfies $\Var(B_i(s,t)-B_i(s',t'))\le C_R(|s-s'|+|t-t'|)$: write the difference as
$(B_i(s,t)-B_i(s',t))+(B_i(s',t)-B_i(s',t'))$ and use $\Var(Y+Z)\le2\Var Y+2\Var Z$, the two terms having variances
$t|s-s'|$ and $s'|t-t'|$. Hence $\E|B(s,t)-B(s',t')|^q\le C_{q,R}\max(|s-s'|,|t-t'|)^{q/2}$ for every $q\ge2$, and the
Kolmogorov continuity theorem for two-parameter processes (see for instance~\cite{Kho02}) gives a version
satisfying~\eqref{eq:holder} for every $\alpha<1/2-2/q$. Taking $q$ large gives every $\alpha<1/2$. The sheet $B$ is
continuous, so it coincides with this version on $R$ almost surely.
\end{proof}

\begin{proof}[Proof of Theorem~\ref{thm:main} for $d=4$]
Fix $K=[a,b]$ with $0<a<b$ and $I=[A_1,A_2]$, $J=[A_3,A_4]$ with $0\le A_1<A_2<A_3<A_4$, all endpoints dyadic rationals.
Fix $\eta\in(0,1]$ and put $\alpha=1/2-\eta/16$, so that $2-4\alpha=\eta/4$. Fix an integer $M\ge1$. For $n$ large
enough that the endpoints lie in $\delta\Z$, where $\delta=2^{-n}$, and that $4M^2\delta^{2\alpha}\le1$, put
$r_n=2M\delta^\alpha$. Then $\delta\le r_n^2\le1$, and
\[
r_n^4\delta^{-2}\bigl(1+\log(1/\delta)\bigr)=16M^4\delta^{-\eta/4}(1+n\log2).
\]
Let $K_\delta=K\cap\delta\Z$, which has at most $C_K\delta^{-1}$ points. By Lemma~\ref{lem:moments}, Markov's
inequality and a union bound, for every integer $p\ge1$,
\[
\PP\Bigl(\max_{s\in K_\delta}N_n(s,r_n)>\delta^{-\eta}\Bigr)\le C_K\delta^{-1}\,\delta^{p\eta}\,C_p\bigl(16M^4\delta^{-\eta/4}
(1+n\log2)\bigr)^p=C_{p,M,K}(1+n)^p\,\delta^{3p\eta/4-1}.
\]
Choose $p$ with $3p\eta/4>1$. The right side is summable in $n$, so by the Borel--Cantelli lemma, almost surely, for all
large $n$,
\begin{equation}\label{eq:allgrid}
N_n(s,r_n)\le\delta^{-\eta}\qquad\text{for every }s\in K_\delta.
\end{equation}
Intersect these events over $M\ge1$, and with the event of Lemma~\ref{lem:holder} for this $\alpha$ and the rectangle
$K\times[0,A_4]$. On the intersection choose an integer $M\ge H$.

Let $s\in K$ be arbitrary and $(u,v)\in Y^*_4(s)\cap(I\times J)$. Round $s$, $u$ and $v$ down to $\delta\Z$, obtaining
$s_n\in K_\delta$, $u_n\in I_\delta$ and $v_n\in J_\delta$; this is possible because the left endpoints lie in
$\delta\Z$. Since $B(s,u)=B(s,v)$, Lemma~\ref{lem:holder} gives
\[
|B(s_n,v_n)-B(s_n,u_n)|\le|B(s_n,v_n)-B(s,v)|+|B(s,u)-B(s_n,u_n)|\le2H\delta^\alpha\le r_n,
\]
so $(u_n,v_n)$ is counted by $N_n(s_n,r_n)$, and $(u,v)$ lies in the square $[u_n,u_n+\delta]\times[v_n,v_n+\delta]$.
By~\eqref{eq:allgrid}, for all large $n$, uniformly in $s\in K$, the set $Y^*_4(s)\cap(I\times J)$ is covered by at most
$\delta^{-\eta}$ squares of side $\delta$. Hence its upper box dimension is at most $\eta$. Intersecting over
$\eta=1/m$, $m\ge1$, gives, almost surely, $\dimB(Y^*_4(s)\cap(I\times J))=0$ for all $s\in K$ simultaneously.

The family of triples $(K,I,J)$ with dyadic endpoints, $a>0$ and $A_2<A_3$, is countable, so almost surely the
conclusion holds for all of them at once. Any compact intervals $I',J'\subset[0,\infty)$ with $I'\cap J'=\emptyset$, say
with $I'$ to the left of $J'$, are contained in dyadic intervals $I\supset I'$ and $J\supset J'$ that are still disjoint,
and every $s>0$ lies in a dyadic $K$ with $a>0$; the case $I'$ to the right of $J'$ follows from the symmetry
$(u,v)\mapsto(v,u)$ of $Y^*_4(s)$. This proves~(b). Every point of $Y^*_4(s)$ lies in one of the countably many sets
$I\times J$ or $J\times I$ with dyadic $I,J$ separated by a positive gap. Since the Hausdorff dimension is at most the
upper box dimension and is stable under countable unions, (a) follows from~(b).
\end{proof}

\begin{proof}[Proof of Theorem~\ref{thm:main} for $d=5$]
The first four coordinates of a five-dimensional Brownian sheet form a four-dimensional Brownian sheet $B'$, and
$Y^*_5(s)\subset\{(u,v):u\ne v,\ B'(s,u)=B'(s,v)\}$ for every $s>0$. The case $d=4$ applied to $B'$ gives the result on
the same event, for all $s>0$ at once.
\end{proof}

\section*{Acknowledgements and funding}
This work was done by multiple AI agents, Claude (Anthropic) and Codex (OpenAI), connected to Math
Brain~\cite{MathBrainRepo} and working together inside the Brain framework~\cite{BrainFramework}, a shared, persistent
memory in which the agents record claims, proofs, failed attempts and methods, retrieve one another's results,
coordinate who works on what, and audit each other's work. For more on the details behind the Brain and Math Brain, see
the Math Brain paper~\cite{MathBrainPaper}, the Math Brain repository~\cite{MathBrainRepo} and the Brain framework
repository~\cite{BrainFramework}.

The research presented in this paper was supported by the European Research Council (ERC) under the European Union's
Horizon Europe research and innovation programme (grant agreement No.~101041711), by the Simons Foundation as part of
the Collaboration on the Mathematical and Scientific Foundations of Deep Learning, by Heights Labs, by Convex Nexus
Capital, by the Israel Science Foundation (grant number 2258/19), and by the Israel Science Foundation (ISF Grant
4101/25). Computational resources were provided in part by Google Cloud through its education credits programme (credit
530550103).

\section*{Declaration of competing interest}
The author holds equity in Heights Labs and in Convex Nexus Capital and has a management role at Convex Nexus Capital.
These entities are included in the funding declaration above.

\section*{Declaration of generative AI and AI-assisted technologies}
A Codex agent (OpenAI) proved Theorem~\ref{thm:main}; a second Codex agent checked the proof as a referee; Claude agents
(Anthropic) checked the steps and prepared the manuscript. The mathematics the paper rests on is due to others. The
problem is Khoshnevisan's~\cite{Kho05}. The existence of exceptional slices and the dimension of $D_d$ are due to
Lyons~\cite{Lyo86} and Mountford~\cite{Mou90}, and the double points of the slices for $d\le3$ to
Penrose~\cite{Pen89,Pen90} and Rosen~\cite{Ros84}. The finiteness of the off-diagonal double-time sets for $d=5$,
which gives both parts of Theorem~\ref{thm:main} in that dimension, is stated by Chen~\cite{Che26}. The method follows
the moment and covering arguments of Khoshnevisan, Wu and Xiao~\cite{KWX06} for the Brownian sheet, with a Brownian
bridge bound for one slice in place of sectorial local nondeterminism. What is new here is Theorem~\ref{thm:main} for
$d=4$, with all slices controlled at once. The author directed the work and is responsible for the
content of the paper and for decisions concerning its dissemination.

\section*{Code availability}
A numerical check of the conditional variance bound of Lemma~\ref{lem:var} on random configurations with repeated
endpoints, and of the lattice sum of Lemma~\ref{lem:lattice}, is in the repository
\begin{center}\url{https://github.com/yspennstate/brownian-sheet-slice-double-times/tree/v1.1}.\end{center}

\appendix
\section{Notation}\label{app:notation}

\begin{center}
\small
\begin{tabular}{lll}
\toprule
Symbol & Meaning & First used\\
\midrule
$B$, $B_i$ & the Brownian sheet and its coordinates & Section~\ref{sec:intro}\\
$D_d$ & the slice parameters whose slice has a double point & Section~\ref{sec:intro}\\
$Y_d(s)$ & Khoshnevisan's double-time set, diagonal included & Section~\ref{sec:intro}\\
$Y^*_d(s)$ & the double times off the diagonal & \eqref{eq:Ystar}\\
$\dimH$, $\dimB$ & Hausdorff and upper box dimension & Section~\ref{sec:intro}\\
$K=[a,b]$ & a compact interval of slice parameters, $a>0$ & Section~\ref{sec:var}\\
$I$, $J$; $A_1,\dots,A_4$ & disjoint compact time intervals and their endpoints & Section~\ref{sec:var}\\
$z_j=(u_j,v_j)$, $X_j$, $X$ & pairs of times and the differences of the slice & Section~\ref{sec:var}\\
$U$, $V$ & the anchors in $I$ and in $J$ & Section~\ref{sec:var}\\
$\rho_U$, $\rho_V$, $\rho$ & distances to the anchors & Section~\ref{sec:var}\\
$W$ & one coordinate of the slice, $t\mapsto B_i(s,t)$ & Lemma~\ref{lem:var}\\
$\mathcal F$, $\mathcal G$ & the $\sigma$-fields of the earlier differences and of the anchors & Lemma~\ref{lem:var}\\
$\sigma^2$ & the conditional variance of one coordinate of $X$ & Lemma~\ref{lem:smallball}\\
$p$ & the order of the moment & Lemma~\ref{lem:moments}\\
$\delta=2^{-n}$; $I_\delta$, $J_\delta$, $K_\delta$ & the grid mesh and the grid points & Section~\ref{sec:moments}\\
$N_n(s,r)$ & the number of near-coincidences & Section~\ref{sec:moments}\\
$\Lambda$ & $r^4\delta^{-2}(1+\log(1/\delta))$ & Lemma~\ref{lem:moments}\\
$H$, $\alpha$ & the H\"older constant and exponent & Lemma~\ref{lem:holder}\\
$\eta$, $M$, $r_n$ & the target dimension, a bound for $H$, the radius $2M\delta^\alpha$ & Section~\ref{sec:proof}\\
\bottomrule
\end{tabular}
\end{center}

\begin{thebibliography}{99}
\bibitem{Che26} K.~Chen, \emph{Uniform dimension results of level-sets of stable sheets}, doctoral thesis no.~11211, EPFL,
Lausanne, 2026. \doi{10.5075/epfl-thesis-11211}
\bibitem{CW25} K.~Chen and G.~Woessner, Zeros of the Brownian sheet, \emph{J.\ Theoret.\ Probab.} \textbf{38} (2025),
article no.~33. \doi{10.1007/s10959-024-01399-1}
\bibitem{DKNWX12} R.~C.~Dalang, D.~Khoshnevisan, E.~Nualart, D.~Wu and Y.~Xiao, Critical Brownian sheet does not have
double points, \emph{Ann.\ Probab.} \textbf{40} (2012), no.~4, 1829--1859. \doi{10.1214/11-AOP665}
\bibitem{DM15} R.~C.~Dalang and C.~Mueller, Multiple points of the Brownian sheet in critical dimensions, \emph{Ann.\
Probab.} \textbf{43} (2015), no.~4, 1577--1593. \doi{10.1214/14-AOP912}
\bibitem{Kho05} D.~Khoshnevisan, Slices of a Brownian sheet: new results and open problems, in \emph{Seminar on
Stochastic Analysis, Random Fields and Applications V}, Progr.\ Probab.\ \textbf{59}, Birkh\"auser, Basel, 2008,
135--174. \doi{10.1007/978-3-7643-8458-6_9}. Equation and section numbers refer to the version
\href{https://arxiv.org/abs/math/0510518}{arXiv:math/0510518}.
\bibitem{Kho02} D.~Khoshnevisan, \emph{Multiparameter Processes: An Introduction to Random Fields}, Springer
Monographs in Mathematics, Springer, New York, 2002. \doi{10.1007/b97363}
\bibitem{KWX06} D.~Khoshnevisan, D.~Wu and Y.~Xiao, Sectorial local non-determinism and the geometry of the Brownian
sheet, \emph{Electron.\ J.\ Probab.} \textbf{11} (2006), paper no.~32, 817--843. \doi{10.1214/EJP.v11-353}
\bibitem{Lyo86} T.~J.~Lyons, The critical dimension at which quasi-every Brownian path is self-avoiding, \emph{Adv.\ in
Appl.\ Probab.} \textbf{18} (1986), suppl.\ \emph{Analytic and Geometric Stochastics}, 87--99.
\bibitem{Mou90} T.~S.~Mountford, Double points and the Ornstein--Uhlenbeck process on Wiener space, \emph{Illinois J.\
Math.} \textbf{34} (1990), no.~1, 38--48. \doi{10.1215/ijm/1255988489}
\bibitem{Pen89} M.~D.~Penrose, On the existence of self-intersections for quasi-every Brownian path in space, \emph{Ann.\
Probab.} \textbf{17} (1989), no.~2, 482--502. \doi{10.1214/aop/1176991411}
\bibitem{Pen90} M.~D.~Penrose, Quasi-everywhere properties of Brownian level sets and multiple points, \emph{Stochastic
Process.\ Appl.} \textbf{36} (1990), no.~1, 33--43. \doi{10.1016/0304-4149(90)90040-Y}
\bibitem{Ros84} J.~Rosen, Self-intersections of random fields, \emph{Ann.\ Probab.} \textbf{12} (1984), no.~1, 108--119.
\doi{10.1214/aop/1176993376}
\bibitem{MathBrainRepo} Y.~Shmalo, \emph{Math Brain: public snapshot, code and tests}, 2026,
\url{https://github.com/yspennstate/math-brain-public}.
\bibitem{MathBrainPaper} Y.~Shmalo, \emph{Math Brain: shared, auditable memory for multi-agent mathematical research},
October 2026, \url{https://github.com/yspennstate/math-brain-public/tree/main/papers/math-brain-2026-10-24}.
\bibitem{BrainFramework} Y.~Shmalo, \emph{Brain Framework}, public software repository, 2026,
\url{https://github.com/yspennstate/brain-framework-public}.
\end{thebibliography}
\end{document}
